\parttitle{II}{Quantum Foundations and the Quantum Reservoir Computing Formulation}

\section{States, observables, and measurement}
\label{sec:quantum-states}

This part fixes the quantum-mechanical machinery in exactly the depth the rest
of the document uses; standard references are
\citet{nielsen2010quantum} for closed-system and circuit material and
\citet{breuer2002theory} for open systems.

A system of $n$ qubits lives on the Hilbert space
$\Hil = (\C^2)^{\otimes n}$ of dimension $d = 2^n$. Pure states are unit
vectors $\ket{\psi} \in \Hil$; the general state is a \emph{density matrix}
$\rho$, a positive semidefinite operator with $\tr\rho = 1$, encompassing pure
states ($\rho = \ketbra{\psi}{\psi}$, $\tr\rho^2 = 1$) and statistical
mixtures ($\tr\rho^2 < 1$). Two structural facts about $\rho$ drive everything
in this document. First, the \emph{exponential dimension}: $\rho$ has
$d^2 - 1 = 4^n - 1$ real parameters, so an $n$-qubit reservoir carries a state
space whose dimension no polynomial-size classical reservoir matches
parameter-for-parameter. Second, the \emph{partial trace}: for a bipartite
system $AB$, the state of $B$ alone is $\rho_B = \tr_A \rho_{AB}$, the unique
operator reproducing all local measurement statistics on $B$. The partial
trace is the mechanism by which the Fujii--Nakajima reservoir
(Section~\ref{sec:qrc-update}) forgets its input qubit while retaining
correlations in the rest --- memory by entanglement.

Observables are Hermitian operators $O$; the Born rule gives the expectation
value $\ev{O} = \tr(O\rho)$. The $n$-qubit Pauli operators
$\{ \id, X, Y, Z \}^{\otimes n}$ form an orthogonal operator basis, so the set
of expectation values of all Pauli strings determines $\rho$ completely; a
reservoir read-out measures a small, fixed subset of them --- in this document
the single-qubit $\ev{Z_i}$, $\ev{X_i}$ and two-qubit $\ev{Z_i Z_j}$. The most
general measurement is a positive operator-valued measure (POVM): a set
$\{E_m\}$ with $E_m \succeq 0$ and $\sum_m E_m = \id$, outcome $m$ occurring
with probability $p_m = \tr(E_m \rho)$. Projective measurement of a qubit in
the computational basis is the special case $E_0 = \ketbra{0}{0}$,
$E_1 = \ketbra{1}{1}$, and is the only primitive current digital hardware
exposes; every ``expectation value'' in an experiment is a finite-sample
average of projective outcomes, a fact whose statistical consequences
(Section~\ref{sec:shot-noise}) shape reservoir design as strongly as any
Hamiltonian choice.

\section{Dynamics: unitaries, channels, and the Lindblad equation}
\label{sec:channels}

\subsection{Closed systems}

A closed system evolves unitarily, $\rho \mapsto U \rho U^\dagger$ with
$U = e^{-iHt}$ generated by a Hamiltonian $H$ (units with $\hbar = 1$
throughout). Unitary evolution is an isometry of state space: it can never
contract two distinct initial states toward each other. An immediate and
frequently overlooked consequence for reservoir computing is that
\emph{a strictly closed quantum system cannot have the echo state property}:
initial-condition differences persist forever, so fading memory
(Definition~\ref{def:fading}) fails. Every working quantum reservoir therefore
contains an explicitly non-unitary element --- an input-injection step that
discards a subsystem, engineered dissipation, or measurement --- and identifying
where that contraction enters is the first question to ask of any proposed
architecture~\citep{chen2019learning,sannia2024dissipation,kobayashi2024coherence}.

\subsection{Quantum channels}

The general physical evolution of a possibly open system is a
\emph{quantum channel}: a completely positive, trace-preserving (CPTP) linear
map $\map{E}$ on density matrices. Every channel admits a Kraus representation
\begin{equation}
  \map{E}(\rho) \;=\; \sum_j K_j \rho K_j^\dagger,
  \qquad
  \sum_j K_j^\dagger K_j = \id ,
  \label{eq:kraus}
\end{equation}
and, equivalently (Stinespring), can be realised as a unitary on the system
plus an environment followed by a partial trace over the environment.
Complete positivity (positivity of $\map{E}\otimes\id$ on any enlarged system)
is what guarantees the map remains physical when the system is entangled with
a bystander; trace preservation is probability conservation. Channels are
exactly the update maps a quantum reservoir may use, and their key analytic
property for this document is \emph{contractivity}: every channel is
non-expansive in trace distance,
$\tfrac12\lVert\map{E}(\rho)-\map{E}(\sigma)\rVert_1 \le
 \tfrac12\lVert\rho-\sigma\rVert_1$,
with strict contraction (for generic state pairs) precisely when the channel
has a non-unitary part. Strict contractivity is the quantum form of the echo
state property (Section~\ref{sec:quantum-esp}).

A convenient linear-algebraic handle used repeatedly in Parts~VII--VIII is the
\emph{Pauli transfer matrix} (PTM): expanding
$\rho = \tfrac1d\sum_\alpha r_\alpha P_\alpha$ over Pauli strings $P_\alpha$
turns a channel into a real matrix $T$ acting on the coefficient vector $r$,
with the unital block's spectral radius governing how fast initial-condition
information decays. \citet{kobayashi2024coherence} sharpen this picture for
reservoirs: fading memory requires the input-independent part of the PTM to be
strictly contractive on the traceless sector, while useful input sensitivity
requires a \emph{coherence influx} --- an input-dependent affine part that keeps
re-exciting the contracting state. Channels that contract without influx (for
example, uniform depolarisation toward the maximally mixed state) satisfy the
ESP but forget the input as fast as the initial condition, and score near
$\NMSE=1$; the distinction is demonstrated numerically in
Part~IV's noise studies.

\subsection{Continuous-time open dynamics}

When the environment coupling is weak and memoryless (Born--Markov), the state
obeys the Gorini--Kossakowski--Sudarshan--Lindblad master
equation~\citep{lindblad1976generators,gorini1976completely}
\begin{equation}
\begin{aligned}
  \dot\rho \;=\; \map{L}(\rho)
  \;&=\; -i[H,\rho]\\
  &\;+\; \sum_k \gamma_k \left( L_k \rho L_k^\dagger
        - \tfrac12\{L_k^\dagger L_k, \rho\} \right),
\end{aligned}
  \label{eq:lindblad}
\end{equation}
whose solution $\rho(t) = e^{\map{L}t}\rho(0)$ is a CPTP semigroup. The jump
operators $L_k$ model the dissipation channels of real hardware: amplitude
damping ($L = \sigma^- = \ketbra{0}{1}$, rate $1/T_1$) relaxes excitation
toward $\ket{0}$; pure dephasing ($L = Z$, rate contribution to $1/T_2$)
destroys phase coherence without moving populations. The distinction is not
cosmetic for reservoirs. Amplitude damping drives the system toward a
\emph{polarised} fixed point ($\ket{0}^{\otimes n}$) from which fresh input
rotations produce large, input-dependent signals, so moderate $T_1$ decay can
supply exactly the contraction the ESP needs while preserving
responsiveness~\citep{sannia2024dissipation,domingo2022optimal}. Strong
dephasing, by contrast, drives states toward the diagonal and, combined with
depolarising elements, toward the maximally mixed state, where all bounded
observables concentrate and the signal is erased --- contraction without
coherence influx~\citep{kobayashi2024coherence,xiong2025role}. ``Noise'' is
therefore neither friend nor enemy of QRC in general; \emph{which} noise, at
what rate, relative to the injection interval, is a design axis
(Part~VIII).

\section{The quantum reservoir computing formulation}
\label{sec:qrc-formulation}

\subsection{State update and read-out}
\label{sec:qrc-update}

A quantum reservoir computer replaces the classical update
Eq.~\eqref{eq:esn-update} by an input-dependent channel acting on the
reservoir's density matrix,
\begin{equation}
\begin{aligned}
  \rho_{t+1}
  \;=\;{}&
  U_{\Delta t}\,
  \Bigl( \ketbra{\psi_{u_{t+1}}}{\psi_{u_{t+1}}} \,\otimes\, \tr_1 \rho_t \Bigr)\,
  U_{\Delta t}^\dagger,\\
  & \ket{\psi_u} = \sqrt{1-u}\,\ket{0} + \sqrt{u}\,\ket{1},
\end{aligned}
  \label{eq:fn-update}
\end{equation}
with $U_{\Delta t} = e^{-iH\Delta t}$ and
\begin{equation}
\begin{aligned}
  H \;=\;{}& \sum_{i<j} J_{ij}\, X_i X_j \;+\; \tfrac{h}{2} \sum_i Z_i ,\\
  & J_{ij} \sim \mathcal{U}[-J/2, J/2].
\end{aligned}
  \label{eq:fn-hamiltonian}
\end{equation}
Equation~\eqref{eq:fn-update} is a legitimate CPTP map (a partial trace, a
tensor product with a fresh state, and a unitary), and each ingredient plays a
distinct role: the partial trace $\tr_1$ is the contraction that grants the
ESP and fading memory; the re-preparation $\ketbra{\psi_u}{\psi_u}$ is the
coherence influx that keeps the reservoir responsive; and the entangling
unitary $U_{\Delta t}$ spreads the fresh input into the retained qubits, where
it interferes with the stored history. Figure~\ref{fig:qrc-pipeline} shows the
full pipeline including measurement and the classical read-out.

\begin{figure*}[t]
\centering
\begin{tikzpicture}[node distance=7mm and 9mm]
  \node[block] (u) {$u_{t+1}$};
  \node[qblock, right=of u] (enc) {encode\\$\ket{\psi_{u_{t+1}}}$ on qubit 1};
  \node[qblock, right=of enc, minimum width=30mm] (evo) {evolve\\$U_{\Delta t}=e^{-iH\Delta t}$};
  \node[qblock, right=of evo] (meas) {measure\\$\ev{O_k}$};
  \node[block, right=of meas] (ro) {ridge\\read-out};
  \node[below=8mm of evo, qblock, minimum width=30mm] (mem) {retained qubits $2..n$:\quad $\tr_1\rho_t$ (memory)};
  \draw[arrow] (u) -- (enc);
  \draw[arrow] (enc) -- (evo);
  \draw[arrow] (evo) -- (meas);
  \draw[arrow] (meas) -- (ro);
  \draw[arrow] (mem.north) -- (evo.south);
  \draw[arrow] (evo.south east) to[out=-45,in=0] node[right,font=\scriptsize]{$t\to t+1$} (mem.east);
  \node[below=1.5mm of ro, font=\scriptsize] {trained};
  \node[above=1.5mm of evo, font=\scriptsize] {fixed};
\end{tikzpicture}
\caption{One step of the Fujii--Nakajima quantum reservoir,
Eq.~\eqref{eq:fn-update}. The input qubit is re-prepared each step (coherence
influx); the remaining qubits carry the fading memory forward; measurement of
a fixed observable set feeds the only trained component, a linear read-out.}
\label{fig:qrc-pipeline}
\end{figure*}

\subsection{Where the nonlinearity enters}
\label{sec:nonlinearity}

Equation~\eqref{eq:qrc-update} is \emph{linear in $\rho$}, and
Eq.~\eqref{eq:qrc-features} is linear too, so a natural worry is that the
whole machine is linear and the read-out learns nothing an autoregression
could not. The worry dissolves once one asks: nonlinear \emph{in what}? Three
distinct mechanisms make the features nonlinear in the \emph{input history},
which is the variable that matters.

\emph{(i) Nonlinear encoding.} The map $u \mapsto \ketbra{\psi_u}{\psi_u}$ in
Eq.~\eqref{eq:fn-update} is quadratic in the amplitudes
$(\sqrt{1-u},\sqrt{u})$, and rotation encodings $R_Y(\gamma u)$ used in
Parts~III--IV make expectation values trigonometric polynomials in $u$. For
parameterised re-uploading circuits this is a theorem: the output is a
truncated Fourier series in the encoded variables, with the accessible
frequency spectrum set by the encoding generators and their
repetitions~\citep{schuld2021effect}. The encoding scale $\gamma$ therefore
controls how much of the read-out's approximation power sits in high
harmonics --- the practical saturation effect quantified in Part~IV.

\emph{(ii) Products through the update composition.} Because the input enters
the channel multiplicatively (the new state is a \emph{product}
$\ketbra{\psi_u}{\psi_u} \otimes \tr_1\rho_t$, then conjugated), composing $T$
steps produces expectation values containing products of functions of
$u_t, u_{t-1}, \dots$ --- exactly the cross terms a NARMA-type target needs.
In the PTM picture: the affine input drive is state-dependent through the
tensor-product structure, so the composed map is polynomial, not affine, in
the input sequence.

\emph{(iii) Measurement-derived nonlinearity.} Any protocol that feeds
measured outcomes back into the drive (Part~III's measurement-based and
feedback architectures~\citep{kobayashi2024feedback,franceschetto2024harnessing}),
or that computes nonlinear functions of estimated expectations classically
(e.g.\ including products $\ev{Z_i}\ev{Z_j}$ as features), imports an
additional, explicitly nonlinear stage. The second-moment observables
$\ev{Z_iZ_j}$ used throughout this document are linear in $\rho$ but carry
information about correlations that no set of single-qubit means determines,
which is why adding them raised the effective feature rank so sharply in the
Part~IV experiments.

What the linearity of Eq.~\eqref{eq:qrc-update} \emph{does} imply is an exact
finite-dimensional bookkeeping: the reservoir state is a vector in
$\R^{4^n-1}$ evolving affinely for each fixed input value, so QRC never
escapes the information-processing capacity bound of
\citet{dambre2012information} --- at most $4^n-1$ independent functionals, and
in practice as many as the \emph{measured, well-conditioned} directions
support. The realistic promise of QRC is therefore not unbounded expressivity
but a large, physically generated feature bank per physical
component~\citep{fujii2017harnessing,mujal2021opportunities}, whose usable
fraction is decided by measurement statistics
(Section~\ref{sec:shot-noise}).

\subsection{Quantum echo state property and fading memory}
\label{sec:quantum-esp}

\begin{definition}[Quantum ESP]
\label{def:qesp}
The reservoir family $\{\map{E}_{\theta,u}\}$ has the echo state property if
for every admissible input sequence and every pair of initial states
$\rho_0, \rho_0'$,
$\lVert \rho_t - \rho_t' \rVert_1 \to 0$ as $t\to\infty$, where
$\rho_t = \map{E}_{\theta,u_t}\circ\cdots\circ\map{E}_{\theta,u_1}(\rho_0)$.
\end{definition}

Sufficient conditions mirror the classical ones with contraction coefficients
replacing spectral radii: if every $\map{E}_{\theta,u}$ is a strict
contraction in trace distance with a uniform coefficient $c<1$, initial-state
information decays as $c^{\,t}$ and the ESP holds together with fading
memory~\citep{chen2019learning,chen2020temporal,martinez2023finite}. For the
Fujii--Nakajima update the contraction is supplied structurally by the partial
trace; for dissipative reservoirs, by the Lindbladian's spectral gap; for
measurement-based schemes, by the state collapse itself. The refinements a
practitioner needs are (a) the \emph{nonstationary/subspace ESP} of
\citet{kobayashi2024extending}, which handles reservoirs where only a
subsystem or a conditioned sector must forget, and (b) the coherence-influx
criterion of \citet{kobayashi2024coherence} separating useful contraction
(toward an input-sensitive, polarised region) from destructive contraction
(toward the maximally mixed state). Both are cheap to check numerically from
the PTM and are part of the validation suite shipped with this document's code.

\subsection{Universality}
\label{sec:qrc-universality}

Universality theorems for QRC parallel Section~\ref{sec:universality}.
\citet{chen2019learning} established that families of dissipative quantum
systems with contractive dynamics and polynomial read-outs realise dense
subsets of fading-memory maps; \citet{chen2020temporal} adapted the
construction to gate-model devices with a hardware demonstration;
\citet{nokkala2021gaussian} proved universality for continuous-variable
Gaussian reservoirs, notable because Gaussian dynamics are efficiently
classically simulable --- universality alone therefore confers no quantum
advantage; and \citet{martinez2023finite} gave a finite-dimensional treatment
identifying which ingredients (contraction, input-dependence, read-out
algebra) are load-bearing. The practical reading is the same as in the
classical case: universality licenses the paradigm, and everything
interesting --- sample efficiency, feature quality per shot, robustness ---
is decided by the concrete architecture, which is Part~III's subject.

\section{Extracting features: multiplexing and measurement back-action}
\label{sec:extraction}

\subsection{Temporal and spatial multiplexing}

The measured feature count $K$, not the Hilbert-space dimension, bounds the
read-out. Two standard tricks enlarge $K$ at fixed hardware.
\emph{Temporal multiplexing}~\citep{fujii2017harnessing} samples the
observables at $V$ intermediate times within each injection interval,
$x_t^{(k,v)} = \tr\bigl(O_k\, e^{-iH v\Delta t/V}\rho_t^{\mathrm{inj}}
e^{+iH v\Delta t/V}\bigr)$, multiplying the feature count by $V$ (the
\emph{virtual nodes}); the document's reference implementation uses $V=4$.
\emph{Spatial multiplexing}~\citep{nakajima2019boosting} runs $M$ independent
small reservoirs on disjoint qubit sets and concatenates their features,
trading entanglement breadth for shot-parallelism and noise isolation ---
usually the right trade on current hardware, and the reason Part~VII's
pipelines shard across small registers rather than maximising register size.

\subsection{Measurement back-action: the four protocols}
\label{sec:backaction}

Equation~\eqref{eq:qrc-features} silently assumes access to exact expectation
values of an undisturbed state, but projective measurement collapses the state
it reads. Four protocols manage the conflict, and choosing among them is a
first-class architectural decision~\citep{mujal2023time}:

\emph{Restart (rewind).} To obtain features at step $t$, re-run the input
sequence from scratch (in practice, from a finite window of length $L$, valid
precisely because fading memory suppresses older inputs) and measure once at
the end. Cost per time step grows with $L\times$ shots, but the state before
measurement is exact. All gate-model examples in Part~IV use this protocol.

\emph{Ensemble.} Run many copies in parallel and measure each copy once per
step --- natural for NMR spin ensembles~\citep{negoro2018machine} and atomic
ensembles, where the ``copies'' are physical molecules or atoms.

\emph{Weak measurement.} Couple a meter weakly, extracting a little
information per step at the cost of a little disturbance;
\citet{mujal2023time} quantify the regime where the running estimate is
accurate while the reservoir dynamics remain effectively intact, and
\citet{franceschetto2024harnessing} show the disturbance itself can be
harnessed as a computational resource.

\emph{Mid-circuit measurement and reset.} Measure a designated subset of
qubits projectively each step and reset it for the next injection --- the
hardware-native realisation of Eq.~\eqref{eq:fn-update}'s partial trace. Its
decisive property, established theoretically and on superconducting hardware
by \citet{hu2024overcoming} (the NISQRC scheme), is that the usable memory
time is set by the \emph{engineered} contraction, not by $T_1/T_2$: the
persistent qubits are refreshed through interaction rather than left to
decohere, so temporal processing can run for durations far exceeding
coherence times. Part~IV demonstrates the primitive in Qiskit with per-step
classical registers.

\subsection{Shot noise}
\label{sec:shot-noise}

With $S$ projective shots, the estimator of a Pauli expectation
$z = \ev{P}$ is a scaled binomial average with standard deviation
\begin{equation}
  \sigma(\hat z) \;=\; \sqrt{\frac{1-z^2}{S}} \;\le\; \frac{1}{\sqrt{S}},
  \label{eq:shot-noise}
\end{equation}
so halving the feature noise costs $4\times$ the shots. Shot noise enters the
ridge problem as noise on the \emph{design matrix}, not the targets, which
biases as well as inflates variance; its interaction with the regulariser is
derived and measured in Part~IV, and its fundamental consequences ---
an effective, noise-dependent cap on the number of usable feature directions
(\emph{eigentasks}) --- are established in
\citet{hu2023tackling} and, for training specifically, in
\citet{ahmed2025optimal}. Two design rules follow and are enforced throughout
this document: always measure the exact-simulation ceiling before any shot
study, so that shot effects are not confused with architectural ones; and
treat $S$ as a first-class resource reported next to every hardware-relevant
number, since a reservoir that needs $10^6$ shots per feature to beat a
\num{30}-neuron ESN has lost on the only axis that counts.

\subsection{Exponential concentration}
\label{sec:concentration}

The scaling threat specific to quantum feature maps is
\emph{concentration}: for deep scrambling dynamics, random circuits, or
noise-dominated evolution, local observables concentrate exponentially in $n$
around state-independent values, so the input-dependent signal shrinks beneath
any fixed shot budget --- the reservoir analogue of barren
plateaus~\citep{mcclean2018barren}. For QRC the phenomenon and its
symmetry-resolved structure are analysed in
\citet{sannia2025exponential,xiong2025role}: scrambling that is useful at
small depth becomes signal-destroying at large depth or size, and noise
accelerates the collapse. The operational defences --- moderate evolution per
step, polarised fixed points, local observables on shallow light-cones,
spatial multiplexing instead of monolithic growth, and empirical monitoring of
feature variance and effective rank as $n$ grows --- recur as design principles
in Part~VIII. Concentration is the single strongest reason this document's
claims are confined to the small-$n$ regime actually simulated, and why no
extrapolation to large registers is offered anywhere.
